\documentclass[11pt]{article}

\usepackage[a4paper,margin=1in]{geometry}
\usepackage{enumitem}
\usepackage{hyperref}
\usepackage[numbers,sort&compress]{natbib}

\hypersetup{
  colorlinks=true,
  linkcolor=black,
  citecolor=black,
  urlcolor=black
}

\setlength{\parindent}{0pt}
\setlength{\parskip}{0.6em}
\setlist[itemize]{leftmargin=1.4em,itemsep=0.2em,topsep=0.2em}
\setlist[enumerate]{leftmargin=1.5em,itemsep=0.2em,topsep=0.2em}

\newcommand{\papertitle}{From greenness to function in ecosystem monitoring}

\begin{document}

\begin{center}
{\Large \textbf{Proposal synopsis}}\\[0.25em]
{\large \textit{Nature Reviews Earth \& Environment}}\\[0.45em]
{\large \textbf{\papertitle}}
\end{center}

\textbf{Article type.} Perspective, approximately 4,000--4,500 words, five display items.

\textbf{Proposed title.} \papertitle.

\textbf{Abstract (about 200 words).}
Ecosystem monitoring is becoming continuous in space and time, but much ecological interpretation still relies on greenness and broad structural proxies. These indicators are indispensable for detecting change, yet they often conflate pigment dynamics, canopy state, water limitation, phenology, disturbance history and viewing geometry. We propose a Perspective on the transition from spectral snapshots to functional ecosystem monitoring. The core argument is that imaging spectroscopy should be treated as a source of plant functional trait anchors, not as a complete monitoring system on its own. Dense multispectral archives can provide temporal continuity, while LiDAR, SAR and InSAR should constrain canopy structure where three-dimensional information is needed. The resulting target is a set of continuous functional observables: spatially and temporally explicit trait estimates with uncertainty, provenance and domain-of-applicability information. The article will distinguish spectrally anchored traits from traits inferred through covariation, temporal context and ecological correlation. It will also show how such observables could support stress-syndrome attribution, resistance, recovery, legacy-effect analysis, dynamic functional diversity and Earth-system model evaluation. A compact European forest proof-of-concept will be used only as an illustrative demonstration, not as a case-study centrepiece.

\textbf{Proposed main sections.}
\begin{enumerate}
\item \textbf{Greenness is not enough.} Why detection-oriented greenness monitoring is insufficient for functional diagnosis under climate extremes.
\item \textbf{Imaging spectroscopy makes function observable, but not continuous.} What spectroscopy can and cannot directly constrain, including the boundary between biochemical trait anchors and structural information.
\item \textbf{From spectral snapshots to continuous functional observables.} A sparse-to-dense framework that links imaging spectroscopy, multispectral time series, structural observations and uncertainty-aware product design.
\item \textbf{What continuous functional observables enable.} Stress syndromes, resistance, recovery, legacy effects, dynamic functional diversity and model benchmarking.
\item \textbf{Guardrails and roadmap.} Reporting standards, domain-of-applicability masks, uncertainty calibration, ecological validation and open spectra-trait-time benchmarks.
\end{enumerate}

\textbf{Display items.}
\begin{enumerate}
\item Conceptual shift from greenness proxies to functional observables.
\item Complementary observing systems and sparse-to-dense monitoring framework.
\item Compact European forest proof-of-concept.
\item Trait anomalies for stress, resistance, recovery and legacy effects.
\item Roadmap and reporting contract for credible functional ecosystem monitoring.
\end{enumerate}

\textbf{Key references and why they matter.}
\begin{itemize}
\item \citet{pettorelli2018satellite}: establishes the ecological importance and limitations of satellite biodiversity and ecosystem monitoring.
\item \citet{zeng2022optical}: summarizes the power and interpretive ambiguity of optical vegetation indices.
\item \citet{cavender2025spectral}: frames spectral biology as a bridge between spectra and biological organization.
\item \citet{cherif2023spectra}: demonstrates transferable multi-trait modelling from heterogeneous spectra-trait data.
\item \citet{mederer2025plant}: shows the importance of curated trait data for hyperspectral trait retrieval.
\item \citet{ji2024transferability}: provides objective evidence that calibration diversity and transferability are central constraints for spectral-trait models.
\item \citet{ji2025transferlearning}: illustrates how transfer learning can broaden leaf trait prediction across space, species and time.
\item \citet{ji2025prisma}: demonstrates spaceborne imaging spectroscopy for seasonal plant trait variability across forest contexts.
\item \citet{dubayah2024structural} and \citet{lang2024gedi}: support the boundary that three-dimensional canopy structure is more directly constrained by active remote sensing than by hyperspectral reflectance alone.
\item \citet{liu2024spectraphenology}: motivates the use of multispectral temporal context for trait and phenology reconstruction.
\end{itemize}

\textbf{Authors and affiliations.}
Chaonan Ji (ENS Paris-Saclay);
Hannes Feilhauer (Leipzig University);
Eya Cherif (Leipzig University);
David Montero (Leipzig University);
Meredith Schumann (University of Zurich);
Julia Joswig (Leipzig University);
Ran Wang (University of Nebraska--Lincoln);
Akpona Okujeni (GFZ German Research Centre for Geosciences);
Gabriele Facciolo (ENS Paris-Saclay);
Philippe Ciais (LSCE/IPSL);
Miguel D. Mahecha (Leipzig University).

\bibliographystyle{unsrtnat}
\bibliography{NREE_proposal_references}

\end{document}
